% Shared building block for slides 34 and 35: eight sub-domains, each a grid
% with a ring of ghost cells, arranged 4x2 with a shared-memory node boundary
% drawn around them. The two slides differ only in what happens *inside* that
% boundary -- a copy, or direct access -- so the layout is written once.

\newcommand{\gsc}{0.155}   % cell size of a sub-domain

% \subdomain{(x,y)} -- a 5x5 interior with an orange ghost ring around it
\newcommand{\subdomain}[1]{%
  \begin{scope}[shift={#1}]
    \foreach \c in {0,...,6}{\foreach \r in {0,...,6}{%
      \ifthenelse{\c=0 \OR \c=6 \OR \r=0 \OR \r=6}
        {\draw[draw=uibkorange!70!black,line width=0.35pt,fill=uibkorange]
           ({\c*\gsc},{-\r*\gsc}) rectangle ++(\gsc,-\gsc);}
        {\draw[draw=blu,line width=0.35pt]
           ({\c*\gsc},{-\r*\gsc}) rectangle ++(\gsc,-\gsc);}}}%
  \end{scope}}

% \ghostgrid{inner arrow style} -- the 4x2 arrangement plus the node boundary
\newcommand{\ghostgrid}[1]{%
  \foreach \k in {0,...,3}{\foreach \m in {0,1}{%
    \subdomain{({\k*1.75},{-\m*1.75})}}}
  % exchanges inside the shared-memory node
  \foreach \k in {0,1,2}{\foreach \m in {0,1}{%
    \draw[<->,#1] ({\k*1.75+1.12},{-\m*1.75-0.54})
               -- ({\k*1.75+1.70},{-\m*1.75-0.54});}}
  \foreach \k in {0,...,3}{%
    \draw[<->,#1] ({\k*1.75+0.54},-1.14) -- ({\k*1.75+0.54},-1.70);}
  % exchanges that must cross the network. A sub-domain is 7*\gsc tall, so the
  % second row ends at -2.835 and the node boundary has to clear it.
  \foreach \k in {0,...,3}{%
    \draw[<->,draw=uibknavy!40!blu,line width=1.4pt]
      ({\k*1.75+0.54},0.22) -- ({\k*1.75+0.54},0.76);
    \draw[<->,draw=uibknavy!40!blu,line width=1.4pt]
      ({\k*1.75+0.54},-3.06) -- ({\k*1.75+0.54},-3.60);}
  \draw[<->,draw=uibknavy!40!blu,line width=1.4pt] (-0.76,-1.42) -- (-0.22,-1.42);
  \draw[<->,draw=uibknavy!40!blu,line width=1.4pt] (6.52,-1.42) -- (7.06,-1.42);
  % the node boundary
  \draw[draw=blu,line width=0.8pt,dash pattern=on 3pt off 3pt]
    (-0.22,0.22) rectangle (6.52,-3.06);}
